\documentclass[10pt,a4paper]{amsart}
\usepackage[margin=19mm]{geometry}
\usepackage{amsmath,amssymb,mathtools}
\usepackage[scr=rsfs,cal=euler]{mathalfa}
\usepackage{xfrac}
\usepackage[unicode,hidelinks,bookmarks=false]{hyperref}
\newcommand{\ie}{i.e.\ }
\newcommand{\cf}{cf.\ }
\newcommand{\resp}{respectively\ }
\newcommand{\Subsection}{\subsection}
\providecommand{\nobreakdash}{\nobreak\hbox{-}}
\setlength{\emergencystretch}{2em}
\multlinegap0pt
\usepackage[shortlabels]{enumitem}
\usepackage{amssymb}
\usepackage{bm}
\usepackage{tensor}
\usepackage{mathtools}
\newtheorem{lem}{Lemma}
\newtheorem{prp}{Proposition}
\newtheorem{thm}{Theorem}
\newcommand{\Gu}{G^{(0)}}
\newcommand{\abs}[1]{\left\vert#1\right\vert}
\newcommand{\cg}{\mathscr{C}_c(G)}
\newcommand{\norm}[1]{\left\|#1\right\|}
\newcommand{\osupp}{\supp^{\circ}}
\newcommand{\supp}{{\rm supp}}
\title{On dense subalgebras of the singular ideal in groupoid C*-algebras}
\author{Julian Gonzales, Jeremy B. Hume}
\date{Source: arXiv:2511.04391v2 (CC BY 4.0)}
\begin{document}
\maketitle

\noindent\emph{Source and licence.} On dense subalgebras of the singular ideal in groupoid C*-algebras. Version arXiv:2511.04391v2 (CC BY 4.0), distributed by arXiv under the Creative Commons Attribution 4.0 International licence, http://creativecommons.org/licenses/by/4.0/, as recorded in the arXiv licence metadata for this exact version and shown on the abstract page https://arxiv.org/abs/2511.04391v2. Full licence notice: \texttt{licenses/arxiv-2511.04391-CC-BY-4.0.txt}.\par\smallskip

\noindent\textit{Editorial scope.} Counted unit: the article's thm SpanJcg (source0.tex lines 900--902). The article's complete original proof of it is delivered with it (source0.tex lines 987--999, one \texttt{proof} environment of 13 lines, which the article places away from the statement). No mathematics is written, repaired or re-derived: the statement and the proof are the article's own lines, in the article's own notation, and no retained line is substituted or re-flowed.  Retained uncounted as the source-local context the proof needs: lines 809--823; lines 835--842; lines 855--855; lines 869--871.  Nothing was omitted from any retained span, so the package carries no omission record.  No retained line contains a citation, so the wrapper carries no bibliography and no bibliography entry.  No retained line contains a figure environment, an image-inclusion command or drawing code, so no image is delivered and the package declares no asset dependency.  Editorial formatting only, in addition: an AMS \texttt{amsart} preamble that restates the article's own declarations and macros used by the retained lines, namely the environments \texttt{lem}, \texttt{prp}, \texttt{thm} on separate counters, together with the standard packages those lines need.  The retained environments print in this document as Lemma 1, Prp 1, Theorem 1 in order of appearance, whereas the article prints the same items with the numbering of its own full document.  The complete native source of the article as distributed in the arXiv e-print for this exact version is bundled unchanged as \texttt{sources/188584/source0.tex}.
\begin{lem}\label{lem:NewWI_condition}
    Let $G$ be an \'etale groupoid that is covered by countably many open bisections.  Fix $x \in \Gu$, and let $g_1, \dots, g_n \in G_{x}$ be distinct points.
    
    There exist open bisections $U_1, \dots, U_n \subseteq G$ satisfying the following conditions:
    \begin{enumerate}[label=(\roman*)]
        \item $g_i \in U_i$ for all $i \in \{1, \dots, n\}$;
        \item $s(U_i) = s(U_j)$ for all $i, j \in \{1, \dots, n\}$;
        \item $W_I = \emptyset$ whenever $I \subseteq \{1, \dotsc, n\}$ is \underline{not} of the form
        \begin{equation}\label{eq:WI_NiceForm}
        \{g_i \colon i \in I\} = kX\cap \{g_1, \dots, g_n\}
        \end{equation}
        for some $k \in G_{x}$ and $X \in \mathcal{X}(x)$.
    \end{enumerate}
    Here $W_I \coloneq s\big( \bigcap_{i \in I} U_i \setminus \bigcup_{j \notin I} U_j \big) \cap C$.    
\end{lem}

\begin{prp}\label{NewCutting}
    Let $G$ be an \'etale groupoid that is covered by countably many open bisections. Fix $x \in \Gu$, and suppose $b \in \mathbb{C}[G^x_x]$ belongs to $\ker(\lambda_{G^x_x/X})$ for all $X \in \mathcal{X}(x)$. Let $\{g_1, \dots, g_n\} \subseteq G^x_x$ denote the support of $b$, let $h \in G_x$, and let $U_1, \dots, U_n$ be open bisections containing $hg_1, \dots, hg_n$ and satisfying the conditions of Lemma~\ref{lem:NewWI_condition}. Then
    \begin{equation}
    \label{f^psi}
    f^\psi \coloneq \sum_{i = 1}^n b(g_i) (\psi \circ s \vert_{U_i})
    \end{equation}
    belongs to $J \cap \cg$ for any $\psi \in C_c\big(s(U_i)\big)$ (note that this does not depend on $i$).
\end{prp}

\subsection{Building blocks for $J \cap \cg$}\label{subsec:BuildingBlocksJcg}

\begin{equation}\label{eq:building_blocks}
f^\Xi \coloneq \sum_{i=1}^n b(g_i) (\psi \circ s \vert_{U_i}).
\end{equation}

\begin{thm}\label{SpanJcg}
    Let $G$ be an \'etale groupoid that is covered by countably many open bisections. Then $\mathcal{E}$ is dense in $J \cap \cg$ with respect to the norm $\norm{\cdot}_I$.
\end{thm}

\begin{proof}
    Given $f \in J \cap \cg$, write $f = \sum_{j=1}^m f_j$ for some $f_j \in C_c(V_j)$ and open bisections $V_j$, and let $H \subseteq G$ be the subgroupoid generated by the open bisections $V_1, \dots, V_m$. Then $H$ is an open subgroupoid that is covered by countably many open bisections. By working in $H$ instead, we can assume that the groupoid $G$ is covered by countably many open bisections. Then, by Theorem~\ref{SpanJcg}, it suffices to assume that $f = f^\Xi$ for some tuple $\Xi = (x, h, b, U_1, \dots, U_n, \psi)$ as in Subsection~\ref{subsec:BuildingBlocksJcg}, and where
    \[
    f^\Xi = \sum_{i=1}^n b(g_i) (\psi \circ s|_{U_i}),
    \]
    as in \eqref{eq:building_blocks}. Take a compact set $K \subseteq G$ containing each $U_i$. As in the proof of Theorem~\ref{SpanJcg} it suffices to find, for each $\varepsilon>0$, a function $f' \in J_{\mathbb{C}}$ satisfying $\osupp(f') \subseteq K$ and $\norm{f^\Xi-f'}_\infty \le \varepsilon$.
    
    Fix $\varepsilon>0$. Define $U \coloneq s(U_i)$ (note that this does not depend on $i$). Since $U$ is totally disconnected, there exists a locally constant function $\varphi \in C_c(U)$ satisfying $\norm{\psi - \varphi}_\infty \le \frac{\varepsilon}{n \max_i(\abs{b(g_i)})}$. Consider the function
    \[
    f' \coloneq \sum_{i=1}^n b(g_i) (\varphi \circ s|_{U_i}).
    \]
    It is clear that $f'$ belongs to the Steinberg algebra $\mathbb{C}G$, and satisfies the conditions $\osupp(f') \subseteq K$ and $\norm{f^\Xi-f'}_\infty \le \varepsilon$. Moreover, by Proposition~\ref{NewCutting}, $f'$ belongs to the singular ideal $J$, and hence also to $J_\mathbb{C}$. This completes the proof.
\end{proof}

\end{document}
